\chapter{Thermal Balance Module}

Thermal managament problems involve the balance between the heat produced and the heat dissipated. Both quantities may be computed by several models. The aim of ThermalBalance (\textit{module\_name} =  \keyword{thermal}) is to treat in a indipendent way the heating and the dissipation to provide users an high flexibility (see Fig.\ref{fig:hmodel}).
\begin{figure}[tbh]
\begin{center}
\includegraphics[width=8cm]{figures/hmodel.eps}
\caption{Heating and dissipation models}\label{fig:hmodel}
\end{center}
\end{figure}

\section{Bulk models}

\subsection{Heating}

Two types of heating models (\keyword{physical\,model heat\_source}) have been implemented. A constant heat source can be inserted with \keyword{type = constant}. The option \keyword{H} sets the value with the units $W/cm^{-3}$.

Example:
\begin{verbatim}
   physical_model heat_source
   {		
     type = constant
     H = 1e7
   }
\end{verbatim}


With \keyword{type = driftdiffusion} we include the heat generated by the transport of electrons ans holes, given by
\begin{equation}\label{Hs}
H =-\nabla \cdot \left[(P_n T + \phi_n)\mathbf{J_n} + (P_p T + \phi_p)\mathbf{J_p}\right] 
\end{equation}
where $\mathbf{J_n}$, $P_n$ and $\phi_n$ are the electrical current, thermoelectric power and electrochemical potential of electrons, respectively. The holes related quantities are p-labeled.
The name of the drift diffusion simulation must be provided (with the options \keyword{drift\_diffusion\_simulation}).

Example:
\begin{verbatim}
   physical_model heat_source
   {		
     model = driftdiffusion

     simulation_name = dd_simul_name
   }
\end{verbatim}

More heating models can be defined in the same regions.

\subsection{Dissipation}
As dissipation models (\keyword{physical\,model heat\_transport} ) \keyword{ThermalBalance} provides the Fourier model and the gray model. The FourierModel (\keyword{type = fourier} ) links the lattice temperature $T$ with the thermal flux $\mathbf{J}$ as
\begin{equation}\label{heat_law}
\mathbf{J} = - \kappa \nabla T  
\end{equation}
where $\kappa$ the thermal conductivity.

The energy conservation is ensured by the continuity equation for the thermal flux, i.e.
\begin{equation}\label{continuity}
\nabla\cdot\mathbf{J} = H
\end{equation}
where $H$ is the total heat source.

Example:
\begin{verbatim}
   physical_model heat_transport
   {		
     type = fourier
   }
\end{verbatim}
The theoretical model of the heat balance problem can be found in Ref.\cite{wachutka}.

The gray model is the simplest approximation of the Boltzmann Transport Equation (BTE) for phonons. It adopts the relaxation time approximation. It further assumes the scattering time to be indipendent on phonon energy and all phonons to have the same group velocity, equal to the first sound velocity.

The gray model (\keyword{type = gray}) is described by the expression
\begin{equation}\label{gray}
  v_g \nabla \cdot \left(\mathbf{s} e''\right)= \frac{e^0- e''}{\tau} + \hat H
\end{equation}
where $e''$ is the density of energy of phonons propagating
toward the $\mathbf s$ direction, $v_g$ is the first sound velocity,
$e^0$ the equilibrium density energy, $\tau$ the scattering relaxation
time and $\hat H$ the total heat source per unit of solid angle.
As Eq.~\ref{gray} has to be solved for each direction, we need to discretize the solid angle. 
According to Fig.\ref{fig:sphere}, the discretization is identified by the number of polar slices $\theta$ (indicated with the option \keyword{N\_theta}) and azimutal slices $\phi$ (indicated with the option \keyword{N\_phi}). The default scheme depends on the dimensionality of the device. 
\begin{figure}[tbh]
\begin{center}
\includegraphics[width=4cm]{figures/sphere.eps}
\caption{Spherical coordinates}\label{fig:sphere}
\end{center}
\end{figure}

The equilibrium energy is given by
\begin{equation}\label{cons}
  e^0 = \frac{1}{4 \pi} \int_{\Omega} e'' d\Omega
\end{equation}
Eq.~(\ref{cons}) is used in this work to achieve energy conservation.

To get the temperature from the density energy we use
\begin{equation}\label{capacity}
  T = T_0 + \frac{4 \pi}{C}e_0
\end{equation}
where the heat capacity, $C$, is considered temperature independent.

The thermal flux can be computed as an angular avarage of the solutions
\begin{equation}\label{flux}
  J^q_i =  \int_{\Omega} v_g e'' s_i d\Omega
\end{equation}

In order to be consistent with the kynetics theory, the relaxation time is computed by the relantioship $\tau = \frac{3 \kappa}{C \,v_g^2}$. The heat capacity (\keyword{C} expressed in $J/K cm^3$, section \keyword{heat\_capacity/constant}) and the sound velocity  (\keyword{vg} expressed in $cm/s$, section \keyword{sound\_velocity/constant}) are taken from the database.
Eqs.~(\ref{cons})-(\ref{gray}) have to be solved in a self-consistent way. The options \keyword{max\_error} and \keyword{max\_iter} set the minimum convergence error and the maximum number of the iteration, respectively.

Sometimes it might be useful to do the domain partition following a criterium bieing beyond the typical geometrical regions. \keyword{ThermalBalance} allows to use a method based on the temperature map computed by the Fourier model. A  threshold in the temperature may act as delimiter between the Fourier and Gray domain. This option can be employed by the option \keyword{partitioning = automatic}. The threshold temperature can be inserted with \keyword{temp\_threshold}.

Exaemple:
\begin{verbatim}
   physical_model heat_transport
   {		
     type = gray
     N_theta = 6
     N_phi = 3
     max_error = 1e3
     max_iter  = 10
     partitioning = automatic
     temp_threshold = 310
   }
\end{verbatim}
the output keyword \keyword{Partitioning} shows how the simulation domain is partitioned into the Fourier and the Gray domain.
 By default we have  \keyword{partitioning = manual} and the region assignement is performed following the usual keyword \keyword{region}.

\subsection{Thermal conductivity}

\keyword{ThermalBalance} implements the constant thermal conductivity model (\keyword{physical\,model thermal\_conductivity},\,\keyword{type = constant}). The three diagonal values of the thermal conductivity tensor (expressed in $W/{cm K}$) can be inserted by the keyword \keyword{ThemCond}. The related database section is \keyword{thermal\_conductivity/constant}.

Example:
\begin{verbatim}
   physical_model thermal_conductivity
   {		
     type = constant
     ThermCond = (1.5,1.5,2.6)
   }
\end{verbatim}

\section{Boundary models}

By default the boundary of the simulation region is thermally insulating, i.e. 
\begin{equation}\label{TH}
\mathbf{J}\cdot \mathbf{N} = 0
\end{equation}
Dirichelt boundary condition, $T = T_{ext}$, can be included with \keyword{type = heat\_reservoir}
\begin{verbatim}
BC_Regions
{
  BC_Region name_BC_region
  {
    type = heat_reservoir
    temperature = 300
  }
}
\end{verbatim}

With \keyword{type = surface\_resistance} it is possible to take into account a thermally resistive interface, i.e.

\begin{equation}\label{TH}
\mathbf{J_q}\cdot \mathbf{N} =  \frac{(T - T_{ext})}{R_s}
\end{equation}

where $R_s$ is the thermal surface resistance and $T_{ext}$ is the external temperature.
Example:

\begin{verbatim}
BC_Regions
{
  BC_Region name_BC_region
  {
    type = surface_resistance
   r_surf = 0.01
    temperature = 300
  }
}
\end{verbatim}


Furthermore, it is possibile to fix the normal thermal power density to a  given value $J_{ext}$, i.e.

\begin{equation}\label{TH}
\mathbf{J_q}\cdot \mathbf{N} = J_{ext}
\end{equation}
This condition is set with \keyword{type = surface\_resistance}.

Example:
\begin{verbatim}
BC_Regions
{
  BC_Region name_BC_region
  {
    type = heat_flux
    power_density = 100 
  }
}
\end{verbatim}
where \keyword{heat\_flux} is expressed in $W/cm^2$.

The thermal boundary conditions are summirezed in Tab.~\ref{table:tbc}:

\begin{table}[htbp]
\center
\begin{tabular}{l||l||l||l}
\textit{description} & \textit{type} & \textit{parameters} & \textit{units}\\
\hline\hline
Ideal insulating interface & (Default) & (No parameters) &   \\
Ideal conducting interface & \texttt{heat\_reservoir} & \texttt{temperature} & $K$\\
Resistive interface &\texttt{surface\_resistance } & \texttt{temperature} & $K$ \\
 & & \texttt{r\_surf} & $cm^2 K W^{-1}$\\
Power density condition& \texttt{heat\_flux} & \texttt{power\_density} & $Wcm^{-2}$
\end{tabular}
\caption{Thermal boundary conditions}
\label{table:tbc}
\end{table}
\section{Output data}
 

\section{Output}

Lattice temperature, total heat source, thermal flux, thermal conductivity and domain partitioning can be plotted by using the plot keyword \keyword{thermal}.
